% @see : https://coopmaths.fr/alea/?uuid=1957a&id=auto6P1A-flash1&alea=e6vK&uuid=adac4&id=auto6P1A-1&n=1&d=10&s=3&s2=5&cd=1&alea=T0on&uuid=adac4&id=auto6P1A-1&n=1&d=10&s=3&s2=1&cd=1&alea=YXJV&uuid=adac4&id=auto6P1A-1&n=1&d=10&s=3&s2=3&cd=1&alea=1fVY&uuid=99d95&id=CM2D1A-1&n=4&d=10&s=false&s2=3&s3=2&s4=false&cd=1&alea=kE0S&pdfParam=eyJ0aXRsZSI6IiIsInJlZmVyZW5jZSI6IiIsInN1YnRpdGxlIjoiIiwic3R5bGUiOiJDb29wbWF0aHMiLCJmb250T3B0aW9uIjoiU3RhbmRhcmRGb250IiwidGFpbGxlRm9udE9wdGlvbiI6MTIsImR5c1RhaWxsZUZvbnRPcHRpb24iOjE0LCJjb3JyZWN0aW9uT3B0aW9uIjoiQXZlY0NvcnJlY3Rpb24iLCJxcmNvZGVPcHRpb24iOiJTYW5zUXJjb2RlIiwidHlwZUZpY2hlIjoiRmljaGUiLCJkdXJhdGlvbkNhbk9wdGlvbiI6IjkgbWluIiwidGl0bGVPcHRpb24iOiJTYW5zVGl0cmUiLCJuYlZlcnNpb25zIjoxLCJleG9zIjp7fX0\%3D&v=eleve
\documentclass[a4paper,11pt,fleqn]{article}

%%% COULEURS %%%
\usepackage[table,svgnames]{xcolor}
\definecolor{nombres}{cmyk}{0,.8,.95,0}
\definecolor{gestion}{cmyk}{.75,1,.11,.12}
\definecolor{gestionbis}{cmyk}{.75,1,.11,.12}
\definecolor{grandeurs}{cmyk}{.02,.44,1,0}
\definecolor{geo}{cmyk}{.62,.1,0,0}
\definecolor{algo}{cmyk}{.69,.02,.36,0}
\definecolor{correction}{cmyk}{.63,.23,.93,.06}
\arrayrulecolor{couleur_theme} % Couleur des filets des tableaux

\usepackage[most]{tcolorbox} % Supprimé à cause de conflits avec ProfCollege
\tcbuselibrary{documentation}

%%%%%%%% Modification paramétrage pour footnotes
% Le paquet semble déjà appelé en amont, curieux qu'il n'y ait pas eu de clash
% avec l'appel ci-dessous mais seulement quand j'ai essayé de l'inclure avec l'option lualatex

% \usepackage{hyperref}
% Parametrages
\hypersetup{
    colorlinks=true,% On active la couleur pour les liens. Couleur par défaut rouge
    linkcolor=black,% On définit la couleur pour les liens internes
    % filecolor=magenta,% On définit la couleur pour les liens vers les fichiers locaux      
    urlcolor=black,% On définit la couleur pour les liens vers des sites web
    % pdftitle={Puissance Quatre},% On définit un titre pour le document pdf
    % pdfpagemode=FullScreen,% On fixe l'affichage par défaut à plein écran
    }
% Pour avoir les numérotations arabic y compris dans les environnements tcolorbox
\renewcommand{\thempfootnote}{\arabic{mpfootnote}}
%%%%%%%%

\usepackage{multicol} 
\usepackage{calc}
\usepackage{enumitem}
\usepackage{graphicx}
\usepackage{tabularx}
\usepackage{tablvar}

\setlength{\parindent}{0mm}
\renewcommand{\arraystretch}{1.5}
\renewcommand{\labelenumi}{\textbf{\theenumi{}.}}
\renewcommand{\labelenumii}{\textbf{\theenumii{}.}}
\setlength{\fboxsep}{3mm}

\setlength{\headheight}{14.5pt}

\setlength{\premulticols}{0pt}

\newcounter{ExoMA}

\newlength{\parindentMA}%
\setlength{\parindentMA}{\parindent}
\addtolength{\parindentMA}{30pt}

\usepackage{simplekv}

\setKVdefault[Theme]{Coopmath,Classiques=false,CANs=false}
\defKV[Theme]{Classique=\setKV[Theme]{Coopmath=false}\setKV[Theme]{Classiques}}
\defKV[Theme]{CAN=\setKV[Theme]{Coopmath=false}\setKV[Theme]{CANs}}

\usepackage{fancyhdr}
\fancyhead[C]{}
\fancyfoot{}

\def\bla{}

\NewDocumentCommand\Theme{o m m m m}{%
  \useKVdefault[Theme]%
  \setKV[Theme]{#1}%
  \ifboolKV[Theme]{CANs}{%
    \usepackage[left=1.5cm,right=1.5cm,top=2cm,bottom=2cm]{geometry}%
    \fancypagestyle{premierePage}{%
      \fancyhead[C]{\textsc{\ifx\bla#2\bla Course aux nombres\else#2\fi}}%
      \fancyfoot[R]{\scriptsize Coopmaths.fr -- CC-BY-SA}%
      \setlength{\headheight}{14.5pt}%
      \NewDocumentCommand\dureeCan{}{\ifx\bla#4\bla 9 minutes\else #4\fi}
      \NewDocumentCommand\titreSujetCan{}{\ifx\bla#3\bla\else\textbf{#3}\fi}
    }%
    \pagestyle{premierePage}
    \colorlet{couleur_theme}{black}%
    \colorlet{couleur_numerotation}{couleur_theme}%
    \let\EXO\EXOCan\let\endEXO\endEXOCan%
  }{%
    \renewcommand{\headrulewidth}{0pt}
    \renewcommand{\footrulewidth}{0pt}
    \pagestyle{fancy}
    \ifboolKV[Theme]{Classiques}{%
      \usepackage[left=0.65cm,right=0.65cm,top=2cm,bottom=1.5cm]{geometry}
      \let\EXO\EXOlibre\let\endEXO\endEXOlibre%
      \colorlet{couleur_theme}{black}%
      \colorlet{couleur_numerotation}{couleur_theme}%
      \fancyhead[C]{\bfseries #3}
      \fancyhead[L]{\bfseries #4}
      \fancyfoot[R]{#5}
    }{%
      \usepackage[left=1.5cm,right=1.5cm,top=4cm,bottom=2cm]{geometry}
      \theme{#2}{#3}{#4}{#5}%
      \let\EXO\EXOcoop\let\endEXO\endEXOcoop%
    }%
  }%
}%

\NewDocumentEnvironment{EXOCan}{m m +b}{%
  #3
}

\NewDocumentEnvironment{EXOcoop}{m m +b}{%
  \begin{tcolorbox}[%
    enhanced,
    breakable,
    colback=white,
    colframe=white,
    coltitle=black,
    title=\IfNoValueTF{#1}{}{\textmd{#1}},%
    overlay unbroken and first={%
      \IfNoValueTF{#2}{}{%
        \node[%
        fill=white,
        anchor=east,
        xshift=10pt,
        text=black
        ]
        at (frame.north east)
        {\scriptsize #2};
      }
      \node[anchor=west,xshift=-20pt,yshift=-15pt] at (frame.north west) {\stepcounter{ExoMA}\numb{\arabic{ExoMA}}};
    }
    ]
    #3
  \end{tcolorbox}
}

\NewDocumentEnvironment{EXOlibre}{m m +b}{%
  \begin{tcolorbox}[%
    enhanced,
    breakable,
    colback=white,
    colframe=white,%couleur_theme,
    coltitle=black,
    title=\stepcounter{ExoMA}\textbf{Exercice \arabic{ExoMA}},%
    ]
    \IfNoValueTF{#1}{}{#1\par}%
    #3
  \end{tcolorbox}%
}


%%% MISE EN PAGE %%%
\usepackage{setspace}
\usepackage{pgf}%
\usepackage{pgfplots}
\pgfplotsset{compat=1.18}
\usetikzlibrary{babel,arrows, arrows.meta,calc,fit,patterns,plotmarks,shapes.geometric,shapes.misc,shapes.symbols,shapes.arrows,shapes.callouts, shapes.multipart, shapes.gates.logic.US,shapes.gates.logic.IEC, er, automata,backgrounds,chains,topaths,trees,petri,mindmap,matrix, calendar, folding,fadings,through,positioning,scopes,decorations.fractals,decorations.shapes,decorations.text,decorations.pathmorphing,decorations.pathreplacing,decorations.footprints,decorations.markings,shadows}

\spaceskip=2\fontdimen2\font plus 3\fontdimen3\font minus3\fontdimen4\font\relax %Pour doubler l'espace entre les mots
\newcommand\numb[1]{ % Dessin autour du numéro d'exercice
  \begin{tikzpicture}[overlay,scale=.8]%,yshift=-.3cm
    \draw[fill=couleur_numerotation,couleur_numerotation](-.4,-0.1)rectangle($(-0.4,.8)+(2em,0)$);
    \draw (-.4,.8) node[white,anchor=north west,inner sep=0.5pt]{\bfseries EX}; 
    \draw[line width=.05cm,couleur_numerotation,fill=white] (0,0)--(.5,.5)--(1,0)--(.5,-.5)--cycle;
    \draw (0.5,0) node[anchor=center,couleur_numerotation]{\large\bfseries#1};
  \end{tikzpicture}
}

% echelle pour le dé
\def\globalscale {0.04}
% abscisse initiale pour les chevrons
\def\xini {3}

\newcommand\theme[4]{%
  \fancyhead[C]{%
    % Tracé du dé
    \begin{tikzpicture}[y=0.80pt, x=0.80pt, yscale=-\globalscale, xscale=\globalscale,remember picture, overlay,transform canvas={xshift=-0.5\linewidth,yshift=2.7cm},fill=couleur_theme]%,xshift=17cm,yshift=9.5cm
      %%%% Arc supérieur gauche%%%%
      \path[fill](523,1424)..controls(474,1413)and(404,1372)..(362,1333)..controls(322,1295)and(313,1272)..(331,1254)..controls(348,1236)and(369,1245)..(410,1283)..controls(458,1328)and(517,1356)..(575,1362)..controls(635,1368)and(646,1375)..(643,1404)..controls(641,1428)and(641,1428)..(596,1430)..controls(571,1431)and(538,1428)..(523,1424)--cycle;
      %%%% Dé face supérieur%%%%
      \path[fill](512,1272)..controls(490,1260)and(195,878)..(195,861)..controls(195,854)and(198,846)..(202,843)..controls(210,838)and(677,772)..(707,772)..controls(720,772)and(737,781)..(753,796)..controls(792,833)and(1057,1179)..(1057,1193)..controls(1057,1200)and(1053,1209)..(1048,1212)..controls(1038,1220)and(590,1283)..(551,1282)..controls(539,1282)and(521,1278)..(512,1272)--cycle;
      %%%% Dé faces gauche et droite%%%%
      \path[fill](1061,1167)..controls(1050,1158)and(978,1068)..(900,967)..controls(792,829)and(756,777)..(753,756)--(748,729)--(724,745)..controls(704,759)and(660,767)..(456,794)..controls(322,813)and(207,825)..(200,822)..controls(193,820)and(187,812)..(187,804)..controls(188,797)and(229,688)..(279,563)..controls(349,390)and(376,331)..(391,320)..controls(406,309)and(462,299)..(649,273)..controls(780,254)and(897,240)..(907,241)..controls(918,243)and(927,249)..(928,256)..controls(930,264)and(912,315)..(889,372)..controls(866,429)and(848,476)..(849,477)..controls(851,479)and(872,432)..(897,373)..controls(936,276)and(942,266)..(960,266)..controls(975,266)and(999,292)..(1089,408)..controls(1281,654)and(1290,666)..(1290,691)..controls(1290,720)and(1104,1175)..(1090,1180)..controls(1085,1182)and (1071,1176)..(1061,1167)--cycle;
      %%%% Arc inférieur bas%%%%
      \path[fill](1329,861)..controls(1316,848)and(1317,844)..(1339,788)..controls(1364,726)and(1367,654)..(1347,591)..controls(1330,539)and(1338,522)..(1375,526)..controls(1395,528)and(1400,533)..(1412,566)..controls(1432,624)and(1426,760)..(1401,821)..controls(1386,861)and(1380,866)..(1361,868)..controls(1348,870)and(1334,866)..(1329,861)--cycle;
      %%%% Arc inférieur gauche%%%%
      \path[fill](196,373)..controls(181,358)and(186,335)..(213,294)..controls(252,237)and(304,190)..(363,161)..controls(435,124)and(472,127)..(472,170)..controls(472,183)and(462,192)..(414,213)..controls(350,243)and(303,283)..(264,343)..controls(239,383)and(216,393)..(196,373)--cycle;
    \end{tikzpicture}
    \begin{tikzpicture}[line cap=round,line join=round,remember picture, overlay, shift={(current page.north west)},yshift=-8.5cm]
      \fill[fill=couleur_theme] (0,5) rectangle (21,6);
      \fill[fill=couleur_theme] (\xini,6)--(\xini+1.5,6)--(\xini+2.5,7)--(\xini+1.5,8)--(\xini,8)--(\xini+1,7)-- cycle;
      \fill[fill=couleur_theme] (\xini+2,6)--(\xini+2.5,6)--(\xini+3.5,7)--(\xini+2.5,8)--(\xini+2,8)--(\xini+3,7)-- cycle;  
      \fill[fill=couleur_theme] (\xini+3,6)--(\xini+3.5,6)--(\xini+4.5,7)--(\xini+3.5,8)--(\xini+3,8)--(\xini+4,7)-- cycle;   
      \node[color=white] at (10.5,5.5) {\LARGE \bfseries{ \MakeUppercase{#4}}};
    \end{tikzpicture}
    \begin{tikzpicture}[remember picture,overlay]
      \node[anchor=north east,inner sep=0pt] at ($(current page.north east)+(0,-.8cm)$) {};
      \node[anchor=west, fill=white, inner sep = 0pt] at ($(current page.north west)+(0.2,-2.3cm)$) {\footnotesize \bfseries{MathALÉA}};
    \end{tikzpicture}
    \begin{tikzpicture}[remember picture,overlay]
      \node[anchor=north east,inner sep=0pt] at ($(current page.north east)+(0,-.8cm)$) {};
      \node[anchor=east, fill=white] at ($(current page.north east)+(-2,-1.5cm)$) {\Huge \textcolor{couleur_theme}{\bfseries{\#}} \bfseries{#2} \textcolor{couleur_theme}{\bfseries \MakeUppercase{#3}}};
    \end{tikzpicture}
  }%
  \fancyfoot[R]{%
    \begin{tikzpicture}[remember picture,overlay]
      \node[anchor=south east] at ($(current page.south east)+(-2,0.25cm)$) {\scriptsize {\bfseries \href{https://coopmaths.fr/}{Coopmaths.fr} -- \href{http://creativecommons.fr/licences/}{CC-BY-SA}}};
    \end{tikzpicture}
    \begin{tikzpicture}[line cap=round,line join=round,remember picture, overlay,xscale=0.5,yscale=0.5, shift={(current page.south west)},xshift=35.7cm,yshift=-6cm]
      \fill[fill=couleur_theme] (\xini,6)--(\xini+1.5,6)--(\xini+2.5,7)--(\xini+1.5,8)--(\xini,8)--(\xini+1,7)-- cycle;
      \fill[fill=couleur_theme] (\xini+2,6)--(\xini+2.5,6)--(\xini+3.5,7)--(\xini+2.5,8)--(\xini+2,8)--(\xini+3,7)-- cycle;  
      \fill[fill=couleur_theme] (\xini+3,6)--(\xini+3.5,6)--(\xini+4.5,7)--(\xini+3.5,8)--(\xini+3,8)--(\xini+4,7)-- cycle;  
    \end{tikzpicture}
  }
  \fancyfoot[C]{}
  \colorlet{couleur_theme}{#1}
  \colorlet{couleur_numerotation}{couleur_theme}
  \def\iconeobjectif{icone-objectif-#1}
  \def\urliconeomethode{icone-methode-#1}
}%

\newcommand*\tikzfootMA{%
  \ifboolKV[Theme]{CANs}{
    \begin{tikzpicture}[remember picture,overlay,shift=(current page.north west)]
      \begin{scope}[x=\textwidth-\oddsidemargin,y=\textheight+5pt]
        \draw[correction,line width=4pt,dashed,dash pattern= on 10pt off 10pt,shift={(-\oddsidemargin,\topmargin)}] (0,0) rectangle (1,-1);
      \end{scope}
    \end{tikzpicture} 
  }{%
  \ifboolKV[Theme]{Classiques}{%
  \begin{tikzpicture}[remember picture,overlay,shift=(current page.south west)]
    \begin{scope}[x={(current page.south east)},y={(current page.north west)}]
      \draw[correction,line width=4pt,dashed,dash pattern= on 10pt off 10pt] ($(0,0)+(5mm,1cm)$) rectangle ($(1,1)+(-5mm,-1.75cm)$);
    \end{scope}
  \end{tikzpicture} 
  }{%
  \begin{tikzpicture}[remember picture,overlay,shift=(current page.south west)]
    \begin{scope}[x={(current page.south east)},y={(current page.north west)}]
      \draw[correction,line width=4pt,dashed,dash pattern= on 10pt off 10pt] ($(0,0)+(5mm,1.5cm)$) rectangle ($(1,1)+(-5mm,-3.75cm)$);
    \end{scope}
  \end{tikzpicture} 
  }%
  }
}

\newenvironment{Correction}{%
  \setcounter{ExoMA}{0}%
  \cfoot{\tikzfootMA}
}{}%

\newcommand\version[1]{
  \fancyhead[R]{
    \begin{tikzpicture}[remember picture,overlay]
    \node[anchor=north east,inner sep=0pt] at ($(current page.north east)+(-.5,-.5cm)$) {\large \textcolor{couleur_theme}{\bfseries V#1}};
    \end{tikzpicture}
  }
}

\usepackage[french]{babel}
% loadPackagesFromContent
\usepackage{tikz}
\newcommand{\e}{\mathrm{e}}
\usepackage{amssymb}
\usepackage{colortbl}
\usetikzlibrary{patterns,patterns.meta}
\usetikzlibrary{arrows.meta}
\usepackage{etoolbox}
\newbool{dys}
\setbool{dys}{false}          
\ifbool{dys}{
% POLICE DYS
% ===== VARIABLE =====
\def\FontChoisie{Fira} % valeurs possibles : Fira, lmodern, tgheros
\newcommand{\choiceFontsDys}[1]{
\ifstrequal{#1}{Fira}{%
  % Fira Sans + Fira Math
  % Description : Fira Sans est une police moderne, et Fira Math est une version mathématique compatible qui maintient le style sans-serif.
  % Utilisation : Fonctionne bien avec XeLaTeX et LuaLaTeX.
  \usepackage{fontspec}
  \setmainfont{Fira Sans}
  \setsansfont{Fira Sans}
  \usepackage{unicode-math}
  \setmathfont{Fira Math}
}{}
\ifstrequal{#1}{lmodern}{%
  % Latin Modern Sans
  % Description : Une version sans-serif de la célèbre police Latin Modern, qui est compatible avec mathastext.
  % Utilisation : Fonctionne avec pdfLaTeX, XeLaTeX et LuaLaTeX.
  \usepackage{lmodern}
  \renewcommand{\familydefault}{\sfdefault}
  \usepackage{mathastext}
}{}
\ifstrequal{#1}{tgheros}{%
  % TeX Gyre Heros + mathastext
  % Description : TeX Gyre Heros est une alternative à Helvetica, et le package mathastext permet d'utiliser la même police pour les mathématiques.
  % Utilisation : Fonctionne avec pdfLaTeX, XeLaTeX, ou LuaLaTeX.
  \usepackage{tgheros}
  \renewcommand{\familydefault}{\sfdefault}
  \usepackage{mathastext}
}{}
}
% Fira ou lmodern ou tgheros
\choiceFontsDys{\FontChoisie}

%\usepackage{unicode-math}
%\usepackage{fontspec}
%\setmainfont{TeX Gyre Schola}
%\setsansfont{TeX Gyre Schola}
%\setmathfont{TeX Gyre Schola Math}
%\setmainfont{OpenDyslexic}[Scale=1.0]
%\setsansfont{OpenDyslexic}[Scale=1.0]
%\setmathfont{OpenDyslexic}[Scale=1.0,range=up/{Latin,latin,num}]
\usepackage[fontsize=14]{scrextend}
\usepackage{setspace}
\setstretch{1.7}
% Une valeur d'environ 1.2 à 1.7 est souvent recommandée. Cela permet d'augmenter l'espace entre les lignes tout en offrant un léger espacement entre les mots.
\setlength{\spaceskip}{1.2em}  
% Une valeur d'environ 1.2em à 1.5em est couramment conseillée. Cela crée un espace plus ample entre les mots, ce qui peut aider à réduire la fatigue visuelle et à améliorer la fluidité de la lecture.
}{
% POLICE STANDARD
\usepackage[T1]{fontenc} 
\usepackage[scaled=1]{helvet}
\usepackage[fontsize=12]{scrextend}
}

\Theme[Coopmaths]{nombres}{}{}{}
\begin{document}

% @see : https://coopmaths.fr/alea?uuid=1957a&id=auto6P1A-flash1&n=1&d=10&alea=e6vK&cd=1&cols=1
\begin{EXO}{}{auto6P1A-flash1}

 Rémi a compté les vêtements dans une armoire. Les effectifs sont représentés sur le diagramme ci-dessous.\\
    \begin{tikzpicture}[baseline,scale = 0.4]

    \tikzset{
      point/.style={
        thick,
        draw,
        cross out,
        inner sep=0pt,
        minimum width=5pt,
        minimum height=5pt,
      },
    }
    \clip (-5,-1) rectangle (17,6);
    	\draw[color ={black},line width = 1.2,-{Stealth[width=4mm]}] (0,0)--(11,0);
	\draw[color ={black},line width = 1.2] (0,0)--(0,5);
	\draw[color ={black},opacity = 0.4, dashed ] (1,0)--(1,5);
	\draw[color ={black},opacity = 0.4, dashed ] (2,0)--(2,5);
	\draw[color ={black},opacity = 0.4, dashed ] (3,0)--(3,5);
	\draw[color ={black},opacity = 0.4, dashed ] (4,0)--(4,5);
	\draw[color ={black},opacity = 0.4, dashed ] (5,0)--(5,5);
	\draw[color ={black},opacity = 0.4, dashed ] (6,0)--(6,5);
	\draw[color ={black},opacity = 0.4, dashed ] (7,0)--(7,5);
	\draw[color ={black},opacity = 0.4, dashed ] (8,0)--(8,5);
	\draw[color ={black},opacity = 0.4, dashed ] (9,0)--(9,5);
	\draw[color ={black},opacity = 0.4, dashed ] (10,0)--(10,5);
	\draw[color ={black},opacity = 0.4, dashed ] (11,0)--(11,5);
	\draw[color ={black},line width = 1.2] (1,-0.13)--(1,0.13);
	\draw[color ={black},line width = 1.2] (2,-0.13)--(2,0.13);
	\draw[color ={black},line width = 1.2] (3,-0.13)--(3,0.13);
	\draw[color ={black},line width = 1.2] (4,-0.13)--(4,0.13);
	\draw[color ={black},line width = 1.2] (5,-0.13)--(5,0.13);
	\draw[color ={black},line width = 1.2] (6,-0.13)--(6,0.13);
	\draw[color ={black},line width = 1.2] (7,-0.13)--(7,0.13);
	\draw[color ={black},line width = 1.2] (8,-0.13)--(8,0.13);
	\draw[color ={black},line width = 1.2] (9,-0.13)--(9,0.13);
	\draw[color ={black},line width = 1.2] (10,-0.13)--(10,0.13);
	\draw[opacity=0.8] (1,-0.4) node[anchor = center, rotate=0] {\scriptsize \color{black}$1$};
	\draw[opacity=0.8] (2,-0.4) node[anchor = center, rotate=0] {\scriptsize \color{black}$2$};
	\draw[opacity=0.8] (3,-0.4) node[anchor = center, rotate=0] {\scriptsize \color{black}$3$};
	\draw[opacity=0.8] (4,-0.4) node[anchor = center, rotate=0] {\scriptsize \color{black}$4$};
	\draw[opacity=0.8] (5,-0.4) node[anchor = center, rotate=0] {\scriptsize \color{black}$5$};
	\draw[opacity=0.8] (6,-0.4) node[anchor = center, rotate=0] {\scriptsize \color{black}$6$};
	\draw[opacity=0.8] (7,-0.4) node[anchor = center, rotate=0] {\scriptsize \color{black}$7$};
	\draw[opacity=0.8] (8,-0.4) node[anchor = center, rotate=0] {\scriptsize \color{black}$8$};
	\draw[opacity=0.8] (9,-0.4) node[anchor = center, rotate=0] {\scriptsize \color{black}$9$};
	\draw[opacity=0.8] (10,-0.4) node[anchor = center, rotate=0] {\scriptsize \color{black}$10$};
	\draw[opacity=0.8] (-2,1.1) node[anchor = center, rotate=0] {\scriptsize \color{black}$\text{chemises}$};
	\draw[opacity=0.8] (-2,2.6) node[anchor = center, rotate=0] {\scriptsize \color{black}$\text{T-shirts}$};
	\draw[opacity=0.8] (-2,4.1) node[anchor = center, rotate=0] {\scriptsize \color{black}$\text{pulls}$};
	\draw [color={black}] (13,0.5) node[anchor = east,scale=1, rotate = 0] {vêtements};
	\draw[color={black},preaction={fill,color = {blue}, opacity = 0.3},pattern color = {black}, pattern = {Lines[angle=45, distance=10pt, line width=0.3pt]}] (0,0.5)--(0,1.5)--(6,1.5)--(6,0.5)--cycle;
\draw [color={black}] (-0.2,1) node[anchor = west,scale=1, rotate = 0] {};
	\draw[color={black},preaction={fill,color = {red}, opacity = 0.3},pattern color = {black}, pattern = {Lines[angle=45, distance=10pt, line width=0.3pt]}] (0,2)--(0,3)--(2,3)--(2,2)--cycle;
\draw [color={black}] (-0.2,2.5) node[anchor = west,scale=1, rotate = 0] {};
	\draw[color={black},preaction={fill,color = {green}, opacity = 0.3},pattern color = {black}, pattern = {Lines[angle=45, distance=10pt, line width=0.3pt]}] (0,3.5)--(0,4.5)--(10,4.5)--(10,3.5)--cycle;
\draw [color={black}] (-0.2,4) node[anchor = west,scale=1, rotate = 0] {};

\end{tikzpicture}\\
 Combien y a-t-il de vêtements en tout ?
\end{EXO}

% @see : https://coopmaths.fr/alea?uuid=adac4&id=auto6P1A-1&n=1&d=10&s=3&s2=5&alea=T0on&cd=1&cols=1
\begin{EXO}{}{auto6P1A-1}

 Dans le parc naturel de Secai, il y a beaucoup d'animaux.\\ Voici un diagramme qui représente les effectifs de quelques espèces.

\medskip
\begin{tikzpicture}[baseline,scale = 0.5]

    \tikzset{
      point/.style={
        thick,
        draw,
        cross out,
        inner sep=0pt,
        minimum width=5pt,
        minimum height=5pt,
      },
    }
    \clip (-9,-5.567727288213004) rectangle (10.5,11.3);
    	\draw [color={black}] (2,-0.2) node[anchor = west,scale=1, rotate = -66] {Gazelles};
	\draw[color ={black}] (2,-0.1)--(2,0.1);
	\draw [color={black}] (4,-0.2) node[anchor = west,scale=1, rotate = -66] {Crocodiles};
	\draw[color ={black}] (4,-0.1)--(4,0.1);
	\draw [color={black}] (6,-0.2) node[anchor = west,scale=1, rotate = -66] {Hyènes};
	\draw[color ={black}] (6,-0.1)--(6,0.1);
	\draw [color={black}] (8,-0.2) node[anchor = west,scale=1, rotate = -66] {Girafes};
	\draw[color ={black}] (8,-0.1)--(8,0.1);
	\draw[color ={black},line width = 1.2] (0,0)--(10,0);
	\draw[color ={black},line width = 1.2,-{Stealth[width=4mm]}] (0,0)--(0,10);
	\draw[color ={black},opacity = 0.4, dashed ] (0,1)--(10,1);
	\draw[color ={black},opacity = 0.4, dashed ] (0,2)--(10,2);
	\draw[color ={black},opacity = 0.4, dashed ] (0,3)--(10,3);
	\draw[color ={black},opacity = 0.4, dashed ] (0,4)--(10,4);
	\draw[color ={black},opacity = 0.4, dashed ] (0,5)--(10,5);
	\draw[color ={black},opacity = 0.4, dashed ] (0,6)--(10,6);
	\draw[color ={black},opacity = 0.4, dashed ] (0,7)--(10,7);
	\draw[color ={black},opacity = 0.4, dashed ] (0,8)--(10,8);
	\draw[color ={black},opacity = 0.4, dashed ] (0,9)--(10,9);
	\draw[color ={black},opacity = 0.4, dashed ] (0,10)--(10,10);
	\draw[color ={black},line width = 1.2] (-0.13,1)--(0.13,1);
	\draw[color ={black},line width = 1.2] (-0.13,2)--(0.13,2);
	\draw[color ={black},line width = 1.2] (-0.13,3)--(0.13,3);
	\draw[color ={black},line width = 1.2] (-0.13,4)--(0.13,4);
	\draw[color ={black},line width = 1.2] (-0.13,5)--(0.13,5);
	\draw[color ={black},line width = 1.2] (-0.13,6)--(0.13,6);
	\draw[color ={black},line width = 1.2] (-0.13,7)--(0.13,7);
	\draw[color ={black},line width = 1.2] (-0.13,8)--(0.13,8);
	\draw[color ={black},line width = 1.2] (-0.13,9)--(0.13,9);
	\draw[opacity=0.8] (-0.75,1.1) node[anchor = center, rotate=0] {\scriptsize \color{black}$10$};
	\draw[opacity=0.8] (-0.75,2.1) node[anchor = center, rotate=0] {\scriptsize \color{black}$20$};
	\draw[opacity=0.8] (-0.75,3.1) node[anchor = center, rotate=0] {\scriptsize \color{black}$30$};
	\draw[opacity=0.8] (-0.75,4.1) node[anchor = center, rotate=0] {\scriptsize \color{black}$40$};
	\draw[opacity=0.8] (-0.75,5.1) node[anchor = center, rotate=0] {\scriptsize \color{black}$50$};
	\draw[opacity=0.8] (-0.75,6.1) node[anchor = center, rotate=0] {\scriptsize \color{black}$60$};
	\draw[opacity=0.8] (-0.75,7.1) node[anchor = center, rotate=0] {\scriptsize \color{black}$70$};
	\draw[opacity=0.8] (-0.75,8.1) node[anchor = center, rotate=0] {\scriptsize \color{black}$80$};
	\draw[opacity=0.8] (-0.75,9.1) node[anchor = center, rotate=0] {\scriptsize \color{black}$90$};
	\draw [color={black}] (0.5,10.5) node[anchor = east,scale=1, rotate = 0] {Nombre d'individus};
	\draw[color={red},opacity = 0.8,fill opacity = 0.4,preaction={fill,color = {red},opacity = 0.4}] (2,3.3000000000000003) circle (0.2);
	\draw[color={red},opacity = 0.8,fill opacity = 0.4,preaction={fill,color = {red},opacity = 0.4}] (4,1.4000000000000001) circle (0.2);
	\draw[color={red},opacity = 0.8,fill opacity = 0.4,preaction={fill,color = {red},opacity = 0.4}] (6,2.7) circle (0.2);
	\draw[color={red},opacity = 0.8,fill opacity = 0.4,preaction={fill,color = {red},opacity = 0.4}] (8,9.3) circle (0.2);
	\draw[color={lightgray}] (2,3.3)--(4,1.4)--(6,2.7)--(8,9.3);

\end{tikzpicture}\\
\textbf {a.}  Quelle est l'espèce la moins nombreuse ? 
    
    	$\square\;$ Gazelles\qquad $\square\;$ Crocodiles\qquad $\square\;$ Hyènes\qquad $\square\;$ Girafes\qquad \\
\\\textbf {b.}  Quelle est l'espèce la plus nombreuse ? 
    
    	$\square\;$ Hyènes\qquad $\square\;$ Girafes\qquad $\square\;$ Gazelles\qquad $\square\;$ Crocodiles\qquad \\
\\\textbf {c.}  L'espèce la plus nombreuse représente ... 
    
    	$\square\;$ Plus de la moitié des animaux\qquad $\square\;$ Moins de la moitié des animaux\qquad $\square\;$ La moitié des animaux\qquad \\

\end{EXO}

% @see : https://coopmaths.fr/alea?uuid=adac4&id=auto6P1A-1&n=1&d=10&s=3&s2=1&alea=YXJV&cd=1&cols=1
\begin{EXO}{}{auto6P1A-1}

 Dans le parc naturel de Cigel, il y a beaucoup d'animaux.\\ Voici un diagramme qui représente les effectifs de quelques espèces.

\medskip
\begin{tikzpicture}[baseline,scale = 0.5]

    \tikzset{
      point/.style={
        thick,
        draw,
        cross out,
        inner sep=0pt,
        minimum width=5pt,
        minimum height=5pt,
      },
    }
    \clip (-6.599376982139838,-6.57121420234855) rectangle (15,6.6);
    	\draw[color={black},fill opacity = 1.1] (0,0) circle (6);
	\draw[color ={black},opacity = 0.8] (0,0.2)--(0,-0.2);\draw[color ={black},opacity = 0.8] (-0.2,0)--(0.2,0);
	\draw  [color={black},preaction={fill,color = {blue},opacity = 0.7},pattern color = {blue}, pattern = crosshatch dots] (-5.984431021743165,0.43195525925500305) -- (0,0) -- (3.6739403974420594e-16,6) arc (90:175.87155963302752:6) ;
	\draw[color={black},preaction={fill,color = {blue}, opacity = 0.7},pattern color = {blue}, pattern = {Lines[angle=45, distance=10pt, line width=0.3pt]}] (7,0)--(8,0)--(8,1)--(7,1)--cycle;
	\draw [color={black}] (8.5,0.5) node[anchor = west,scale=1.5, rotate = 0] {Servals};
	\draw  [color={black},preaction={fill,color = {GreenYellow},opacity = 0.7},pattern color = {GreenYellow}, pattern = {Lines[angle=0, distance=10pt, line width=0.3pt]}] (-4.2119605197347205,-4.2731005815679115) -- (0,0) -- (-5.984431021743165,0.43195525925500433) arc (175.87:225.41128440366973:6) ;
	\draw[color={black},preaction={fill,color = {GreenYellow}, opacity = 0.7},pattern color = {GreenYellow}, pattern = {Lines[angle=45, distance=10pt, line width=0.3pt]}] (7,1.5)--(8,1.5)--(8,2.5)--(7,2.5)--cycle;
	\draw [color={black}] (8.5,2) node[anchor = west,scale=1.5, rotate = 0] {Guépards};
	\draw  [color={black},preaction={fill,color = {brown},opacity = 0.7},pattern color = {brown}, pattern = dots] (5.924771266515311,-0.9471459441261234) -- (0,0) -- (-4.21196051973472,-4.273100581567913) arc (-134.59:-9.085412844036696:6) ;
	\draw[color={black},preaction={fill,color = {brown}, opacity = 0.7},pattern color = {brown}, pattern = {Lines[angle=45, distance=10pt, line width=0.3pt]}] (7,3)--(8,3)--(8,4)--(7,4)--cycle;
	\draw [color={black}] (8.5,3.5) node[anchor = west,scale=1.5, rotate = 0] {Zèbres};
	\draw  [color={black},preaction={fill,color = {LightSlateBlue},opacity = 0.7},pattern color = {LightSlateBlue}, pattern = grid] (5.551115123125783e-16,6) -- (0,0) -- (5.92477126651531,-0.9471459441261261) arc (-9.08:90.00256880733946:6) ;
	\draw[color={black},preaction={fill,color = {LightSlateBlue}, opacity = 0.7},pattern color = {LightSlateBlue}, pattern = {Lines[angle=45, distance=10pt, line width=0.3pt]}] (7,4.5)--(8,4.5)--(8,5.5)--(7,5.5)--cycle;
	\draw [color={black}] (8.5,5) node[anchor = west,scale=1.5, rotate = 0] {Lycaons};

\end{tikzpicture}\\
\textbf {a.}  Quelle est l'espèce la moins nombreuse ? 
    
    	$\square\;$ Servals\qquad $\square\;$ Lycaons\qquad $\square\;$ Zèbres\qquad $\square\;$ Guépards\qquad \\
\\\textbf {b.}  Quelle est l'espèce la plus nombreuse ? 
    
    	$\square\;$ Guépards\qquad $\square\;$ Lycaons\qquad $\square\;$ Servals\qquad $\square\;$ Zèbres\qquad \\
\\\textbf {c.}  L'espèce la plus nombreuse représente ... 
    
    	$\square\;$ Moins de la moitié des animaux\qquad $\square\;$ Plus de la moitié des animaux\qquad $\square\;$ La moitié des animaux\qquad \\

\end{EXO}

% @see : https://coopmaths.fr/alea?uuid=adac4&id=auto6P1A-1&n=1&d=10&s=3&s2=3&alea=1fVY&cd=1&cols=1
\begin{EXO}{}{auto6P1A-1}

 Dans le parc naturel de Barmwich, il y a beaucoup d'animaux.\\ Voici un diagramme qui représente les effectifs de quelques espèces.

\medskip
\begin{tikzpicture}[baseline,scale = 0.5]

    \tikzset{
      point/.style={
        thick,
        draw,
        cross out,
        inner sep=0pt,
        minimum width=5pt,
        minimum height=5pt,
      },
    }
    \clip (-9,-4.197409101749103) rectangle (10.5,6.3);
    	\draw[color={black},preaction={fill,color = {blue}, opacity = 0.3},pattern color = {black}, pattern = {Lines[angle=45, distance=10pt, line width=0.3pt]}] (1.7,0)--(1.7,3.8)--(2.3,3.8)--(2.3,0)--cycle;
\draw [color={black}] (2,-0.2) node[anchor = west,scale=1, rotate = -66] {Lycaons};
	\draw[color={black},preaction={fill,color = {GreenYellow}, opacity = 0.3},pattern color = {black}, pattern = {Lines[angle=45, distance=10pt, line width=0.3pt]}] (3.7,0)--(3.7,2.3)--(4.3,2.3)--(4.3,0)--cycle;
\draw [color={black}] (4,-0.2) node[anchor = west,scale=1, rotate = -66] {Girafes};
	\draw[color={black},preaction={fill,color = {brown}, opacity = 0.3},pattern color = {black}, pattern = {Lines[angle=45, distance=10pt, line width=0.3pt]}] (5.7,0)--(5.7,0.9)--(6.3,0.9)--(6.3,0)--cycle;
\draw [color={black}] (6,-0.2) node[anchor = west,scale=1, rotate = -66] {Servals};
	\draw[color={black},preaction={fill,color = {LightSlateBlue}, opacity = 0.3},pattern color = {black}, pattern = {Lines[angle=45, distance=10pt, line width=0.3pt]}] (7.7,0)--(7.7,1.5)--(8.3,1.5)--(8.3,0)--cycle;
\draw [color={black}] (8,-0.2) node[anchor = west,scale=1, rotate = -66] {Hyènes};
	\draw[color ={black},line width = 1.2] (0,0)--(10,0);
	\draw[color ={black},line width = 1.2,-{Stealth[width=4mm]}] (0,0)--(0,5);
	\draw[color ={black},opacity = 0.4, dashed ] (0,1)--(10,1);
	\draw[color ={black},opacity = 0.4, dashed ] (0,2)--(10,2);
	\draw[color ={black},opacity = 0.4, dashed ] (0,3)--(10,3);
	\draw[color ={black},opacity = 0.4, dashed ] (0,4)--(10,4);
	\draw[color ={black},opacity = 0.4, dashed ] (0,5)--(10,5);
	\draw[color ={black},line width = 1.2] (-0.13,1)--(0.13,1);
	\draw[color ={black},line width = 1.2] (-0.13,2)--(0.13,2);
	\draw[color ={black},line width = 1.2] (-0.13,3)--(0.13,3);
	\draw[color ={black},line width = 1.2] (-0.13,4)--(0.13,4);
	\draw[opacity=0.8] (-0.75,1.1) node[anchor = center, rotate=0] {\scriptsize \color{black}$10$};
	\draw[opacity=0.8] (-0.75,2.1) node[anchor = center, rotate=0] {\scriptsize \color{black}$20$};
	\draw[opacity=0.8] (-0.75,3.1) node[anchor = center, rotate=0] {\scriptsize \color{black}$30$};
	\draw[opacity=0.8] (-0.75,4.1) node[anchor = center, rotate=0] {\scriptsize \color{black}$40$};
	\draw [color={black}] (0.5,5.5) node[anchor = east,scale=1, rotate = 0] {Nombre d'individus};

\end{tikzpicture}\\
\textbf {a.}  Quelle est l'espèce la moins nombreuse ? 
    
    	$\square\;$ Girafes\qquad $\square\;$ Hyènes\qquad $\square\;$ Servals\qquad $\square\;$ Lycaons\qquad \\
\\\textbf {b.}  Quelle est l'espèce la plus nombreuse ? 
    
    	$\square\;$ Lycaons\qquad $\square\;$ Hyènes\qquad $\square\;$ Servals\qquad $\square\;$ Girafes\qquad \\
\\\textbf {c.}  L'espèce la plus nombreuse représente ... 
    
    	$\square\;$ Moins de la moitié des animaux\qquad $\square\;$ La moitié des animaux\qquad $\square\;$ Plus de la moitié des animaux\qquad \\

\end{EXO}

% @see : https://coopmaths.fr/alea?uuid=99d95&id=CM2D1A-1&n=4&d=10&s=false&s2=3&s3=2&s4=false&alea=kE0S&cd=1&cols=1
\begin{EXO}{Répondre aux questions à l'aide du texte.}{CM2D1A-1}

 Plusieurs amis reviennent du marché. Il s'agit de Albert et Océane.\\Albert rapporte $2$ kiwis, $5$ pêches et $7$ coings. \\Océane rapporte $4$ kiwis et $1$ coing. 

\medskip
\textbf {a.}  Remplir le tableau suivant. \\\textbf {b.}  Quel est le nombre total de fruits achetés par les amis ?\\\textbf {c.}  Qui a rapporté le plus de fruits ?\\\textbf {d.}  Quel fruit a été rapporté en la plus grosse quantité ?\\$\renewcommand{\arraystretch}{1}
\begin{array}{|c|c|c|c|c|}
\hline
\cellcolor{lightgray} \text{Amis\textbackslash fruits} & \cellcolor{lightgray} \text{Coing} & \cellcolor{lightgray} \text{Kiwi} & \cellcolor{lightgray} \text{Pêche} & \cellcolor{lightgray} \text{TOTAL}\\
\hline
\cellcolor{lightgray} \text{Albert} &  &  &  & \\
\hline
 \cellcolor{lightgray} \text{Océane} &  &  &  & \\
\hline
 \cellcolor{lightgray} \text{TOTAL} &  &  &  & \\
\hline
 \end{array}
\renewcommand{\arraystretch}{1}$

\end{EXO}


\clearpage

\begin{Correction}
\begin{EXO}{}{}

 Il y a $6+2+10 = 18$ vêtements en tout.
\end{EXO}

\begin{EXO}{}{}

 \textbf {a.}  L'animal le moins nombreux parmi ces espèces est :  
    
    	$\square\;$ Gazelles\qquad $\blacksquare\;$ Crocodiles\qquad $\square\;$ Hyènes\qquad $\square\;$ Girafes\qquad \\
\\\textbf {b.}  L'animal le plus nombreux parmi ces espèces est :  
    
    	$\square\;$ Hyènes\qquad $\blacksquare\;$ Girafes\qquad $\square\;$ Gazelles\qquad $\square\;$ Crocodiles\qquad \\
\\\textbf {c.}  L'animal le plus nombreux parmi ces espèces représente :  
    
    	$\blacksquare\;$ Plus de la moitié des animaux\qquad $\square\;$ Moins de la moitié des animaux\qquad $\square\;$ La moitié des animaux\qquad \\

\end{EXO}

\begin{EXO}{}{}

 \textbf {a.}  L'animal le moins nombreux parmi ces espèces est :  
    
    	$\square\;$ Servals\qquad $\square\;$ Lycaons\qquad $\square\;$ Zèbres\qquad $\blacksquare\;$ Guépards\qquad \\
\\\textbf {b.}  L'animal le plus nombreux parmi ces espèces est :  
    
    	$\square\;$ Guépards\qquad $\square\;$ Lycaons\qquad $\square\;$ Servals\qquad $\blacksquare\;$ Zèbres\qquad \\
\\\textbf {c.}  L'animal le plus nombreux parmi ces espèces représente :  
    
    	$\blacksquare\;$ Moins de la moitié des animaux\qquad $\square\;$ Plus de la moitié des animaux\qquad $\square\;$ La moitié des animaux\qquad \\

\end{EXO}

\begin{EXO}{}{}

 \textbf {a.}  L'animal le moins nombreux parmi ces espèces est :  
    
    	$\square\;$ Girafes\qquad $\square\;$ Hyènes\qquad $\blacksquare\;$ Servals\qquad $\square\;$ Lycaons\qquad \\
\\\textbf {b.}  L'animal le plus nombreux parmi ces espèces est :  
    
    	$\blacksquare\;$ Lycaons\qquad $\square\;$ Hyènes\qquad $\square\;$ Servals\qquad $\square\;$ Girafes\qquad \\
\\\textbf {c.}  L'animal le plus nombreux parmi ces espèces représente :  
    
    	$\blacksquare\;$ Moins de la moitié des animaux\qquad $\square\;$ La moitié des animaux\qquad $\square\;$ Plus de la moitié des animaux\qquad \\

\end{EXO}

\begin{EXO}{}{}

 \textbf {a.}  Voici le tableau complet. \\$\renewcommand{\arraystretch}{1}
\begin{array}{|c|c|c|c|c|}
\hline
\cellcolor{lightgray} \text{Amis\textbackslash fruits} & \cellcolor{lightgray} \text{Coing} & \cellcolor{lightgray} \text{Kiwi} & \cellcolor{lightgray} \text{Pêche} & \cellcolor{lightgray} \text{TOTAL}\\
\hline
\cellcolor{lightgray} \text{Albert} & 7 & 2 & 5 & 14\\
\hline
 \cellcolor{lightgray} \text{Océane} & 1 & 4 & 0 & 5\\
\hline
 \cellcolor{lightgray} \text{TOTAL} & 8 & 6 & 5 & 19\\
\hline
 \end{array}
\renewcommand{\arraystretch}{1}$
\\\textbf {b.}  Le nombre total de fruits est : $19$.\\\textbf {c.}  On regarde la dernière colonne du tableau. La personne qui a rapporté le plus de fruits est Albert. Ce nombre maximal de fruits est de $14$.\\\textbf {d.}  On regarde la dernière ligne du tableau. Il y a plus de coings que d'autres fruits. Il y en a $8$.
\end{EXO}

\clearpage
\end{Correction}
\end{document}

% Local Variables:
% TeX-engine: luatex
% End: